Examen général de Dominique LOYER
Question 1 :
Il ne vous a pas échappé que votre thématique de recherche a quelque chose à voir avec la
notion de responsabilité. Hans Jonas l’avait dressé au rang de principe. Pour lui, l’espèce
humaine était la mieux dotée au sein de la biosphère pour endosser l’habit de « l’espèce
responsable du devenir du vivant ».
Il ne vous a guère échappé non plus que le numérique vient titiller cette façon d’envisager
l’ordre du monde puisqu’il arrive parfois que l’on demande à un non-humain de rendre des
comptes sur ce qu’il fait et communique. Ceci s’est déjà présenté dans l’histoire de
l’humanité puisque des « animaux » et des « machines » ont été appelés à la barre de
tribunal en qualité d’accusés ou de témoins afin de réguler les affaires humaines.
Dans le 1er paragraphe de votre résumer de sujet de recherche, vous faîtes écho à cette idée
en écrivant : « il devient impossible pour un individu de mettre en doute une décision
automatique prise à son encontre en demandant à l’algorithme de personnalisation de lui
rendre des comptes sur le pourquoi de sa décision ».
Or, là, vous mettez le doigt sur un aspect de votre travail que j’aimerai que vous exploriez
un peu plus : le pouvoir. Max Weber, dans son ouvrage intitulé la bureaucratie, décrit
comment la bureaucratie s’appuie sûr et déploie des machines à gouverner : « La
bureaucratie concentre le pouvoir aux dépens du monde extérieur qui subit l’opacité de
ses règles » (Wikipédia, Bureaucratie, 2025). C’est pourquoi je vous poserai la question
suivante : « Quel lien pourriez-vous faire entre les modalités de fonctionnement de
l’algorithme de personnalisation et le phénomène bureaucratique ? ».

Réf. The Concept of Bureaucracy by Max Weber International Journal of Social
Science Studies.
Vol. 7, No. 2; March 2019
ISSN 2324-8033 E-ISSN 2324-8041
Published by Redfame Publishing
URL: http://ijsss.redfame.com

Réf. David Graeber, Bureaucratie. L'utopie des règles, Paris, Les Liens qui libèrent,
2015, 304 p., traduit de l'anglais par Françoise Chemla, ISBN : 979-10-209-0291-7.


Question 2 : Imaginons : Au fil de votre doctorat, vous mettez à l’épreuve un ou des
algorithmes de personnalisation et vous démontrez des relations de cause à effet dans la
génération de logiques prédatrices pour les usagers et les collectivités.
o Quelles seront les étapes que vous envisageriez de mettre en place pour modifier
les modalités de fonctionnement des algorithmes que vous aurez testés ?

Réf. https://declarationmontreal-iaresponsable.com


Question 3 : Il est probable qu’au cours la démarche développée pour répondre à la
question 2, vous soyez amenés à expérimenter des rapports de force entre des humains à
propos du fonctionnement d’artefacts. Vous considérez que c’est le moment d’en profiter
pour rappeler à l’ensemble de vos interlocuteurs la responsabilité humaine,
organisationnelle et sociale de leurs routines professionnelles.
o En vous inspirant du cas réel de Biosphere Code, comment vous vous y prendriez
pour formuler un code de déontologie des codeurs qui travaillent sur la formulation
d’algorithme de personnalisation ?
Réf.
Réf. https://darkmachines.earth/blog/the-biosphere-code
Réf. https://www.lepoint.fr/editos-du-point/laurence-neuer/algorithmes-ce-sont-les-
individus-qu-il-faut-responsabiliser-29-01-2018-2190382_56.php#11